\section{Orientations and protected integration from the volumes: attack--heal}
\label{sec:orientations-protected-integration-from-volumes-attack-heal}

\providecommand{\Z}{\mathbb Z}
\providecommand{\C}{\mathbb C}
\providecommand{\La}{\Lambda}
\providecommand{\RPerf}{\operatorname{RPerf}}
\providecommand{\RHom}{\operatorname{RHom}}
\providecommand{\sdet}{\operatorname{sdet}}
\providecommand{\CoHA}{\operatorname{CoHA}}
\providecommand{\Hall}{\mathcal H}
\providecommand{\Prim}{\operatorname{Prim}}
\providecommand{\Ran}{\operatorname{Ran}}
\providecommand{\hCS}{\mathrm{hCS}}
\providecommand{\BV}{\mathrm{BV}}
\providecommand{\TS}{\mathrm{TS}}
\providecommand{\nuDelta}{\nu_{\Delta_5}}
\providecommand{\epsor}{\epsilon_o}
\providecommand{\autcor}{\mathfrak g_{\Delta_5}}

\subsection{Claim attacked}

The attacked upgrade is:
\[
\hbox{cross-volume BV/CoHA/orientation technology}
\quad\Longrightarrow\quad
\epsor=\nuDelta .
\]
This is not proved in the present constellation.  The healed statement is
type-correct and conditional:

\begin{quote}
If one constructs an oriented compact \(K3\times E\) Hall/CoHA
realization, a protected integration map, and a lifted action of the
type-II Weyl group on the oriented completed Hall object, then one may ask
whether the induced orientation monodromy character equals the
restriction of the Maass character of \(\Delta_5\).  None of the existing
volumes constructs this comparison character.
\end{quote}

\subsection{Domains of the two characters}

\paragraph{Automorphic character.}
The theorem-level character is
\[
\nuDelta:\Sp_4(\Z)\longrightarrow\{\pm1\}.
\]
It is defined by Maass's multiplier on the standard generators and
satisfies
\[
\Delta_5|_5g=\nuDelta(g)\Delta_5,
\qquad
\Delta_5^2=\Delta_{10}\ \hbox{scalar}.
\]
Local anchors:
\[
\texttt{main.tex:651--683},\qquad
\texttt{main.tex:719--731}.
\]
Via the exterior-square model and the type-II chamber decomposition it
restricts to
\[
\nuDelta|_{W^{(2)}(\La^{2,1}_{II})}:W^{(2)}(\La^{2,1}_{II})\to\{\pm1\},
\qquad
\nuDelta(s_{\delta_i})=-1,
\]
and on this subgroup
\[
\Delta_5(w\cdot Z)=\det(w)\Delta_5(Z).
\]
Local anchors:
\[
\texttt{main.tex:1438--1455},\qquad
\texttt{main.tex:1524--1538}.
\]
Vol III gives the same character by a Steinberg-generator, Weil, and
orthogonal-determinant comparison:
\[
\texttt{calabi-yau-quantum-groups/chapters/examples/cy\_d\_kappa\_stratification.tex:3750--3784},
\]
and records the type-I/type-II sign asymmetry at
\[
\texttt{calabi-yau-quantum-groups/chapters/examples/cy\_d\_kappa\_stratification.tex:4061--4085}.
\]

\paragraph{Geometric orientation character.}
The geometric character is not defined until the following data exist:
\[
\bigl(\Hall^{\mathrm{or}}_{K3\times E},o,\Theta_{\hCS\to\Hall}^{\mathrm{or}},
I^{\mathrm{prot}},\rho_o\bigr),
\]
where \(o\) is a strong square-root orientation system, and
\[
\rho_o:W^{(2)}(\La^{2,1}_{II})
\longrightarrow
\operatorname{Aut}(\Hall^{\mathrm{or}}_{K3\times E})
\]
is a lifted Weyl action preserving the Hall product, completion,
orientation line system, protected trace, and charge-to-root map.  Then
\[
\epsor(w)\in\{\pm1\}
\]
is the sign by which \(w\) transports the chosen orientation branch:
\[
\rho_o(w)^*\ell_\gamma\simeq \epsor(w)\,\ell_{w\gamma}.
\]
Without \(\rho_o\), the expression \(\epsor(w)\) has no domain.
Without \(o\), it has no line on which to act.

The current Igusa manuscript states this as a criterion and open problem:
\[
\texttt{main.tex:903--930},\qquad
\texttt{main.tex:933--953},\qquad
\texttt{main.tex:2179--2235}.
\]

\subsection{Type-correct comparison map}

The comparison is not between a scalar partition function and a sign.
It is an equality of two characters on the same group after two maps have
been constructed:
\[
\begin{tikzcd}[column sep=large]
W^{(2)}(\La^{2,1}_{II})
\arrow[r,"\iota_{\mathrm{ext}}"]
\arrow[d,"\rho_o"']
&
\Sp_4(\Z)/\{\pm I\}
\arrow[d,"\nuDelta"]
\\
\operatorname{Aut}(\Hall^{\mathrm{or}}_{K3\times E})/
\operatorname{Aut}^{+}(\Hall^{\mathrm{or}}_{K3\times E})
\arrow[r,"\operatorname{sgn}_o"']
&
\{\pm1\}.
\end{tikzcd}
\]
Here \(\iota_{\mathrm{ext}}\) is the exterior-square/lattice
identification used in the manuscript, and
\(\operatorname{sgn}_o\circ\rho_o=\epsor\).  The desired assertion is
\[
\operatorname{sgn}_o\circ\rho_o
=
\nuDelta\circ\iota_{\mathrm{ext}}
\quad\hbox{on }W^{(2)}(\La^{2,1}_{II}).
\]
This is meaningful only after the Hall-side action \(\rho_o\) is
constructed.  The automorphic side already exists.

\subsection{BV, CoHA, and orientation anchors}

\begin{center}
\small
\begin{tabular}{p{0.33\textwidth}p{0.58\textwidth}}
\hline
Anchor & Content \\
\hline
\texttt{main.tex:903--930} &
Protected Hall integration is a criterion: compact Hall/CoHA sector,
oriented extension correspondences, vanishing cycles, HN completions,
protected integration, primitive projection, and lifted Weyl action are
all required. \\
\texttt{main.tex:933--953} &
The orientation monodromy character is open; the expected Pfaffian local
model has not been constructed for \(K3\times E\). \\
\texttt{main.tex:2179--2235} &
Compact Igusa datum requires \(E_3\), specialization, oriented Hall
factorization, orientation square roots, Thom--Sebastiani transport, and
the character condition \(\epsor=\nuDelta\) on the type-II Weyl group. \\
\texttt{calabi-yau-quantum-groups/chapters/theory/cy3\_chain\_level\_bridge.tex:522--652} &
Vol III defines an oriented hCS--Hall comparison datum on the DWR/Ran
nerve: BV observables, oriented critical CoHA, determinant-line square
roots, Hall pull-push, and Thom--Sebastiani transport. \\
\texttt{calabi-yau-quantum-groups/chapters/theory/cy3\_chain\_level\_bridge.tex:654--700} &
The DWR/Ran theorem is a sufficiency theorem.  It constructs no datum; it
starts after orientation, trace, twist, completion, product, and
coherence have been supplied. \\
\texttt{calabi-yau-quantum-groups/chapters/theory/cy3\_chain\_level\_bridge.tex:1099--1211} &
The hCS--Hall descent obstruction is
\((o_{\mathrm{MC}},o_{\mathrm{or}},o_{\mathrm{gr}},
o_{\mathrm{TS}},o_{\mathrm{fact}})\); vanishing is necessary and
sufficient for a global oriented comparison on a fixed cover. \\
\texttt{calabi-yau-quantum-groups/chapters/theory/cy3\_chain\_level\_bridge.tex:3452--3488} &
Vol III proves a Hall-side orientation torsor trivialization on a Kummer
DWR cover, but explicitly says this does not kill the relative
hCS--Hall orientation obstruction. \\
\texttt{calabi-yau-quantum-groups/chapters/theory/cy3\_chain\_level\_bridge.tex:3554--3695} &
Finite DWR/Ran comparison exists if finite primitive packages
\((\eta_{\mathrm{MC}},\lambda_{\mathrm{or}},\eta_{\mathrm{gr}},
H_{\mathrm{TS}},H_{\mathrm{fact}},Q)\) are supplied; compact completion
still requires inverse-limit compatibility. \\
\texttt{calabi-yau-quantum-groups/chapters/examples/k3\_chiral\_bialgebra\_platonic.tex:4130--4192} &
Joyce/Kinjo--Park--Safronov orientation data gives generic CY\(_3\)
orientability of the derived moduli stack.  This is orientation
existence, not Weyl reflection monodromy. \\
\texttt{chiral-bar-cobar/worldview\_synthesis\_2026\_04\_17.tex:647--654} &
Vol I supplies the operadic/Ran/BV/BRST heptagon language; no
\(K3\times E\) determinant-line monodromy character is constructed. \\
\texttt{chiral-bar-cobar-vol2/chapters/connections/ht\_bulk\_boundary\_line\_core.tex:1030--1049} &
Vol II proves a boundary Thom--Sebastiani multiplicativity theorem.  It
is not an orientation monodromy computation for the compact Hall theory. \\
\texttt{topological-strings/main.tex:2219--2248} &
The BCOV volume imports Costello's perturbative BV package and the
Costello--Li flat theorem; it does not specialize them to the compact
Igusa Hall orientation problem. \\
\texttt{topological-strings/main.tex:2258--2315} &
The relevant Hamiltonian specialization is itself formulated as an open
BV graph problem. \\
\hline
\end{tabular}
\end{center}

\subsection{Theorem, proposed, open}

\begin{center}
\small
\begin{tabular}{p{0.34\textwidth}p{0.14\textwidth}p{0.42\textwidth}}
\hline
Claim & Status & Reason \\
\hline
\(\nuDelta:\Sp_4(\Z)\to\{\pm1\}\) and its type-II restriction &
Theorem &
Maass character plus exterior-square/reflection calculation. \\
\(\Delta_5(s_{\delta_i}Z)=-\Delta_5(Z)\) &
Theorem &
Simple type-II real-root reflections have Maass value \(-1\). \\
Reduced OP/DT scalar branch
\(-4096\Delta_5^{-2}\) &
Theorem/imported &
Oberdieck--Pixton/Oberdieck--Pandharipande/Oberdieck scalar theory plus
the manuscript's \(D_5=64^{-1}\Delta_5\) normalization. \\
Generic CY\(_3\) orientation data on \(\RPerf(X)\) &
Theorem elsewhere / imported &
KPS/BBDJS orientability supplies square-root orientation data on
shifted-symplectic derived stacks. \\
Hall-side orientation torsor on a Kummer DWR cover &
Proved in Vol III, local target only &
It prepares the Hall target orientation line; Vol III says the relative
hCS--Hall orientation obstruction can still survive. \\
Finite oriented hCS--Hall comparison &
Conditional theorem &
Vol III proves the finite descent statement after the finite primitive
package is supplied. \\
Compact protected integration for \(K3\times E\) &
Open &
The compact extension-closed Hall/CoHA atlas, reduced quotient descent,
HN completion, primitive extraction, and protected trace are not all
constructed. \\
\(\epsor=\nuDelta\) &
Open &
No existing volume constructs \(\rho_o\) and computes determinant-line
monodromy around the type-II reflections. \\
\hline
\end{tabular}
\end{center}

\subsection{Exact missing lemmas}

\begin{enumerate}
\item \textbf{Compact Hall atlas lemma.}  Construct an extension-closed
compact critical Hall/CoHA object for the relevant \(K3\times E\) sector,
before quotienting by \(E\), with finite-type HN truncations and support
property.
\item \textbf{Reduced quotient descent lemma.}  Prove that vanishing-cycle
coefficients, extension correspondences, and orientation local systems
descend through the reduced \(E\)-quotient without losing the relative
position data needed for Hall multiplication.
\item \textbf{Strong orientation lemma.}  Construct square roots
\[
\ell_\gamma^{\otimes2}\simeq
\sdet\RHom(\mathcal E_\gamma,\mathcal E_\gamma)
\]
with direct-sum, exact-triangle, overlap, and triple-extension
Thom--Sebastiani coherences.
\item \textbf{Protected integration lemma.}  Construct
\[
I^{\mathrm{prot}}_{\sigma,o}:
\Hall^{\mathrm{or}}_{K3\times E}\to
\widehat{\C[\Gamma_{\mathrm{eff}}]}
\]
compatible with Hall product, wall crossing, HN completion, and primitive
projection.
\item \textbf{Weyl lift lemma.}  Construct a lifted action
\[
\rho_o:W^{(2)}(\La^{2,1}_{II})
\to\operatorname{Aut}(\Hall^{\mathrm{or}}_{K3\times E})
\]
preserving the charge map
\[
\alpha(n,l,m)=2nf_2-lf_3+2mf_{-2},
\]
the oriented Hall product, and the protected integration map.
\item \textbf{Reflection-normal-mode lemma.}  Identify the wall crossed by
each \(s_{\delta_i}\) with the diagonal real-root normal mode, and prove
that transporting the Pfaffian/determinant square root reverses exactly
one oriented normal factor:
\[
\epsor(s_{\delta_i})=-1.
\]
\item \textbf{Relation lemma.}  Prove that the signs computed on simple
type-II reflections satisfy the Coxeter/Nikulin relations, so they define
a character on all of \(W^{(2)}(\La^{2,1}_{II})\).
\item \textbf{Exterior-square compatibility lemma.}  Prove that the
Hall-side lifted reflection and the automorphic reflection are the same
element under the exterior-square map
\[
W^{(2)}(\La^{2,1}_{II})\to \Sp_4(\Z)/\{\pm I\},
\]
with no hidden central-lift sign.
\item \textbf{Pro-limit coherence lemma.}  Promote the finite DWR/Ran
comparison primitives \((\eta_{\mathrm{MC}},\lambda_{\mathrm{or}},
\eta_{\mathrm{gr}},H_{\mathrm{TS}},H_{\mathrm{fact}},Q)\) through the
\((N,r,L,m)\)-inverse systems; equivalently kill the displayed
\(\varprojlim^1\) obstruction in Vol III.
\end{enumerate}

\subsection{Compact recognition criterion}

Let \(\mathfrak K_X\) be the compact Igusa datum of the manuscript.
The orientation character is recognized as the automorphic character if
the following finite checklist is satisfied.

\begin{enumerate}
\item The oriented compact Hall/CoHA object and protected integration map
exist.
\item The type-II Weyl group acts on the oriented completed Hall object
and preserves the charge-to-root map.
\item The induced orientation transport is multiplicative, hence a
character
\[
\epsor:W^{(2)}(\La^{2,1}_{II})\to\{\pm1\}.
\]
\item For the three simple type-II walls
\[
\delta_1=2f_2-f_3,\qquad
\delta_2=2f_{-2}-f_3,\qquad
\delta_3=f_3,
\]
the determinant-line square root changes sign:
\[
\epsor(s_{\delta_1})=\epsor(s_{\delta_2})=\epsor(s_{\delta_3})=-1.
\]
\item The lift to \(\Sp_4(\Z)/\{\pm I\}\) is the exterior-square lift
used by the automorphic calculation.
\end{enumerate}

Then, because \(W^{(2)}(\La^{2,1}_{II})\) is generated by the type-II
reflections and because the Maass restriction is \(\det\) on that Weyl
group,
\[
\epsor=\nuDelta|_{W^{(2)}(\La^{2,1}_{II})}.
\]
This criterion is compact, but it is not yet satisfied in the volumes.

\subsection{Falsified normalizations and false shortcuts}

\begin{center}
\small
\begin{tabular}{p{0.29\textwidth}p{0.60\textwidth}}
\hline
Shortcut & Failure \\
\hline
\([q^{1/2}r^{1/2}s^{1/2}]\Delta_5=64\) fixes orientation &
False.  The factor \(64\) is the theta-leading coefficient; rescaling
\(\Delta_5\) does not change its character. \\
\(D_5=64^{-1}\Delta_5\) fixes orientation &
False.  This chooses the monic Borcherds product; it constructs no
determinant-line square root on the Hall side. \\
\(-4096\Delta_5^{-2}\) fixes orientation &
False.  The number \(4096=64^2\) and the minus sign are scalar OP
normalization data; the square has forgotten \(\nuDelta\). \\
\(\Delta_{10}=\Delta_5^2\) determines \(\nuDelta\) geometrically &
False for this problem.  \(\Delta_{10}\) is scalar on \(\Sp_4(\Z)\);
the square cannot recover the missing sign without an independent
orientation-monodromy lift. \\
Canonical orientability of \(\RPerf(X)\) gives \(\epsor=\nuDelta\) &
False.  Orientability supplies a square-root system; equality with the
Maass character requires a Weyl action and monodromy computation. \\
Hall-side orientation torsor triviality gives \(\epsor=\nuDelta\) &
False.  Vol III explicitly leaves the relative hCS--Hall orientation
class \(o_{\mathrm{or}}\) and the compact inverse-limit compatibility as
separate data. \\
\(\nuDelta=(-1)^F\) in a super-VOA settles the comparison &
Type error.  \((-1)^F\) is state parity; \(\nuDelta\) is a
modular/Weyl-group character.  A group action on the oriented compact
Hall object must relate them. \\
Diagonal restriction \(Z=\tau I\) makes \(1/\Delta_5\) an ordinary
chiral character &
Insufficient.  It changes the function's domain but does not construct
compact Hall products, orientation-line monodromy, or the type-II Weyl
action. \\
\hline
\end{tabular}
\end{center}

\subsection{Conclusion}

The existing volumes contain substantial orientation technology:
Joyce/KPS orientation data, oriented critical CoHA language,
DWR/Ran hCS--Hall comparison criteria, boundary Thom--Sebastiani
multiplicativity, and Costello BV/factorization machinery.  They do not
construct the missing object:
\[
\rho_o:W^{(2)}(\La^{2,1}_{II})
\to\operatorname{Aut}(\Hall^{\mathrm{or}}_{K3\times E})
\]
nor compute the induced sign on the orientation square root.  Therefore
\[
\epsor=\nuDelta
\]
remains an open orientation-monodromy recognition problem, not a theorem
and not a normalization convention.

\subsection{Checks}

Files read for this report include \(\texttt{CLAUDE.md}\),
\(\texttt{AGENTS.md}\), \(\texttt{main.tex}\), \(\texttt{proj.bib}\),
the local \(\texttt{agent\_material/12\_*}\) and
\(\texttt{agent\_material/13\_*}\) reports relevant to protected
integration, Ran locality, primitive recognition, and synthesis, plus
the cited cross-volume anchors in
\(\texttt{/Users/raeez/calabi-yau-quantum-groups}\),
\(\texttt{/Users/raeez/chiral-bar-cobar}\),
\(\texttt{/Users/raeez/chiral-bar-cobar-vol2}\), and
\(\texttt{/Users/raeez/topological-strings}\).  No main manuscript file
was edited.  No build was run because this standalone report is not
included by \(\texttt{main.tex}\).
